\documentclass[10pt,a4paper]{amsart}
\usepackage[margin=19mm]{geometry}
\usepackage{amsmath,amssymb,mathtools}
\usepackage[scr=rsfs,cal=euler]{mathalfa}
\usepackage{xfrac}
\usepackage[unicode,hidelinks,bookmarks=false]{hyperref}
\newcommand{\ie}{i.e.\ }
\newcommand{\cf}{cf.\ }
\newcommand{\resp}{respectively\ }
\newcommand{\Subsection}{\subsection}
\providecommand{\nobreakdash}{\nobreak\hbox{-}}
\setlength{\emergencystretch}{2em}
\multlinegap0pt
\usepackage{bm}
\usepackage{bbm}
\usepackage{dsfont}
\usepackage{xspace}
\usepackage{cancel}
\usepackage{mathtools}
\usepackage{amsthm}
\makeatletter
\@ifundefined{lemma}{\newtheorem{lemma}{Lemma}}{}
\makeatother
\title[An etude on a renormalization]{An etude on a renormalization}
\author{A. V. Ivanov}
\date{Source: arXiv:2504.12818v1 (CC BY 4.0)}
\begin{document}
\maketitle

\noindent\emph{Source and licence.} An etude on a renormalization. Version arXiv:2504.12818v1 (CC BY 4.0), distributed under the Creative Commons Attribution 4.0 International licence, as recorded in the arXiv licence metadata for this exact version and shown on the abstract page https://arxiv.org/abs/2504.12818v1. Full rights notice: licenses/arxiv-2504.12818-CC-BY-4.0.txt.\par\smallskip

\noindent\textit{Editorial scope.} Counted unit: the Lemma whose source label is \texttt{34-l-4} in the article (source0.tex lines 319--321, source label \texttt{34-l-4}), delivered together with the article's own complete original proof of it (source0.tex lines 322--360, one \texttt{proof} environment of 39 lines). No mathematics is written, repaired or re-derived: statement and proof are the article's own text line for line. Required context retained uncounted: 34-1 (lines 115--119); 34-5 (lines 177--180); 34-l-2 (lines 196--204); 34-17 (lines 312--316). Every definition, display and label of those lines is retained unchanged. Omitted from this package and not redistributed: 1--114, 120--176, 181--195, 205--311, 317--318, 361--755. Nothing inside a retained excerpt is omitted or altered: every retained body is delivered exactly as the article's lines are written, so no omission receipt of any kind accompanies this package. Citations: the retained excerpts carry no citation command at all, so no bibliography entry of the article is transcribed and no reference is redistributed. Reference adaptation: none was needed and none was made; every reference inside the delivered unit resolves inside it, so no unresolved reference or citation is produced. Printed slips kept literal: no printed word, symbol, hypothesis or number of the counted statement or of its proof was altered. Layout and class: the article's own class is \texttt{article} with packages including amsfonts, amssymb, amsmath, amsthm, cite, ucs, inputenc, graphicx, babel, slashed, textcomp, bm, titlesec, tikz-cd; the wrapper is the lane's pinned \texttt{amsart} editorial wrapper at 10pt on A4, which restates the theorem-like environments the retained text needs and adds the article's own macro definitions that the retained text needs (none), with extra wrapper packages (bm, bbm, dsfont, xspace, cancel, mathtools). The delivered numbering is the unit's own. Assets: no retained span contains a figure environment, a graphics-inclusion command, a PDF-page inclusion, a table or a TikZ picture; no image file is copied into this package, the package declares no asset dependency and no image file is redistributed with it. Licence: version arXiv:2504.12818v1 of this article is distributed under the Creative Commons Attribution 4.0 International licence, as recorded in the arXiv licence metadata for this exact version and shown on the abstract page https://arxiv.org/abs/2504.12818v1. A full rights notice is delivered at licenses/arxiv-2504.12818-CC-BY-4.0.txt. Characters: every retained excerpt is pure ASCII, so the delivered engine, XeLaTeX with its default Latin Modern text font, reports no missing character.
\begin{equation}\label{34-1}
\Phi_n(s,\beta)=
\prod_{j=1}^{n}\int_{\mathbb{R}}\frac{\mathrm{d}a_j}{\sqrt{\pi/\beta_j}}
\,e^{-\beta_j^{\phantom{1}}a_j^2+isa_j^2}.
\end{equation}

\begin{equation}\label{34-5}
\Phi_n(s,\beta)=\prod_{j=1}^{n}\Big(1-is/\beta_j\Big)^{-\frac{1}{2}}=
f_n(s)e^{ig_n(s)/2},
\end{equation}

\begin{lemma}\label{34-l-2} Let $s\in\mathbb{R}\setminus\{0\}$ and $\beta\in\mathcal{B}_2$. Then, taking into account all of the above, the following statements are true:
	\begin{enumerate}
		\item $0<f(s)<f_{n+1}(s)< f_{n}(s)<1$ for all $n\in\mathbb{N}$;
		\item $f(\cdot)\in L^p(\mathbb{R})$ for all $p>0$;
		\item $g(s)<+\infty$ if and only if $b_1(\beta)<+\infty$;
		\item $\Phi_n(s,\beta)$ has a limit when $n\to+\infty$ if and only if $b_1(\beta)<+\infty$.
	\end{enumerate}
If $s=0$, then we have the relations $\Phi(0,\beta)=1$, $f(0)=1$, and $g(0)=0$.
\end{lemma}

\begin{equation}\label{34-17}
	Z_n(\lambda,\beta)=
	\Bigg(\prod_{j=1}^{n}\int_{\mathbb{R}}\frac{\mathrm{d}a_j}{\sqrt{\pi/\beta_j}}\Bigg)
	\exp\bigg(-S_0[p_n(a)]-\lambda\Big(S_1[p_n(a)]\Big)^2\bigg).
\end{equation}

\begin{lemma}\label{34-l-4}
Let $\beta\in\mathcal{B}_2\setminus\mathcal{B}_1$, then the function $Z(\lambda,\beta)=0$ for all $\lambda>0$. At the same time $Z(1,\beta)=1$.
\end{lemma}

\begin{proof} Note that the functions from \eqref{34-1} and \eqref{34-17} are related by the following integral transformation
\begin{equation}\label{34-20}
	Z_n(\lambda,\beta)=\mathrm{T}\big(\Phi_n(\,\cdot\,,\beta)\big)(\lambda)=
	\frac{1}{\sqrt{4\pi\lambda}}
	\int_{\mathbb{R}}\mathrm{d}s\,e^{-s^2/(4\lambda)}
	\Phi_n(s,\beta),
\end{equation}
which follows from the calculation of the standard Gaussian integral. Define a set of numbers and a function
\begin{equation}\label{34-23}
c_n=\sum_{j=1}^n\beta_j^{-1},\,\,\,h_n(s)=\frac{1}{\sqrt{4\pi\lambda}}e^{-s^2/(4\lambda)}
f_n(s)e^{i(g_n(s)-sc_n)/2}.
\end{equation}
Due to exponential decrease and properties from Lemma \ref{34-l-2}, the function $h_n(s)$ and any of its derivatives belong to $L^1(\mathbb{R})$. Besides, the estimates are correct
\begin{equation}\label{34-25}
\bigg|\frac{\mathrm{d}f_n(s)}{\mathrm{d}s}\bigg|\leqslant\frac{|s|b_2(\beta)}{2},\,\,\,
\bigg|\frac{1}{2}\frac{\mathrm{d}(g_n(s)-sc_n)}{\mathrm{d}s}\bigg|\leqslant
\sum_{j=1}^n\frac{s^2}{2\beta_j^3}\leqslant\frac{s^2b_2(\beta)}{2\mu(\beta)},
\end{equation}
\begin{equation}\label{34-24}
\bigg|\frac{\mathrm{d}h_n(s)}{\mathrm{d}s}\bigg|\leqslant
\frac{e^{-s^2/(4\lambda)}}{\sqrt{4\pi\lambda}}
\bigg(\frac{|s|}{2\lambda}+\frac{|s|b_2(\beta)}{2}+\frac{s^2b_2(\beta)}{2\mu(\beta)}\bigg).
\end{equation}
Note that the right-hand sides do not depend on the parameter $n$, so the inequalities are also valid for the limiting case.
Next, substituting the representation from \eqref{34-5} into \eqref{34-20} and integrating it by parts, we get
\begin{equation}\label{34-26}
Z_n(\lambda,\beta)=\int_{\mathbb{R}}\mathrm{d}s\,e^{isc_n/2}h_n(s)=-
\frac{2}{ic_n}\int_{\mathbb{R}}\mathrm{d}s\,e^{isc_n/2}\frac{\mathrm{d}h_n(s)}{\mathrm{d}s}.
\end{equation}
Therefore, passing to the absolute values of the functions, we obtain an estimate of the form
\begin{equation}\label{34-18}
|Z_n(\lambda,\beta)|\leqslant
\frac{2}{c_n}
\Bigg(
\int_{\mathbb{R}}\mathrm{d}s\,\frac{e^{-s^2/(4\lambda)}}{\sqrt{4\pi\lambda}}
\bigg(\frac{|s|}{2\lambda}+\frac{|s|b_2(\beta)}{2}+\frac{s^2b_2(\beta)}{2\mu(\beta)}\bigg)\Bigg),
\end{equation}
in which the right-hand side tends to zero at $n\to+\infty$ for all fixed $\lambda>0$. The final equality $Z(1,\beta)=1$ follows from direct substitution.
\end{proof}

\end{document}
